% ECONOMICS_PROBLEM: {"id": "J64-349", "title": "Job-search information and persistent mismatch", "jel_code": "J64", "source_id": "OP-0352"}
% !TeX program = xelatex
\documentclass[11pt,a4paper]{article}
\usepackage[margin=23mm]{geometry}
\usepackage{fontspec}
\setmainfont{Latin Modern Roman}
\usepackage{amsmath,amssymb,mathtools}
\usepackage{xcolor,enumitem,booktabs,longtable,array,xurl,hyperref}
\definecolor{muted}{HTML}{53636C}
\hypersetup{colorlinks=true,urlcolor=blue,linkcolor=blue}
\setlength{\parindent}{0pt}
\setlength{\parskip}{5pt}
\newcommand{\R}{\mathbb R}
\newcommand{\E}{\mathbb E}
\newcommand{\N}{\mathbb N}
\newcommand{\ind}{\mathbf 1}
\newcommand{\Var}{\operatorname{Var}}
\newcommand{\Cov}{\operatorname{Cov}}
\newcommand{\supp}{\operatorname{supp}}
\DeclareMathOperator*{\argmin}{arg\,min}
\DeclareMathOperator*{\argmax}{arg\,max}
\newcommand{\entrymeta}[3]{{\sffamily\small\color{muted}\textbf{\texttt{#1}}\quad|\quad #2\par #3\par}}
\newcommand{\Problemref}[1]{\texttt{#1}}
\newcommand{\JELref}[2]{\texttt{#2} (JEL \texttt{#1})}
\begin{document}
\section*{Job-search information and persistent mismatch}
\textbf{Problem ID:} \texttt{J64-349}\quad \textbf{JEL:} \texttt{J64}\par
% BEGIN ECONOMICS STATEMENT
\label{problem:OP-0352}
\label{audited:ECON-095}
\label{audited:ECON-096}
{\sffamily\small\color{muted}Original subject group: Labor markets and work\par}
\label{definitions:problem:30007634}
\textbf{Audit classification:} Published research agenda. Full manuscript checked. Nonpay uncertainty and social ties are stated directions; the source finds preferences important and does not establish pay misinformation as the dominant explanation.
\subsection*{Definitions and precise research target}
\textbf{Complete editorial problem.} Consider workers searching within a specified labor market. For worker \(i\) and employer \(j\), let \(b_{ij}\) be the worker's belief about pay and nonpay job characteristics, \(v_{ij}\) the value of a fully informed feasible offer, and \(k_{ij}\) the switching cost. An information intervention \(z\in\{0,1\}\) supplies independently validated employer information; it does not guarantee an offer. Let \(J_i(z)\) be the eventual employer and \(Y_{iT}(z)\) discounted earnings through a fixed horizon \(T\). Define \(\tau_T=E[Y_{iT}(1)-Y_{iT}(0)]\) and \(\tau_J=\Pr(J_i(1)\ne J_i(0))\). Randomization identifies marginal employer distributions and \(\tau_T\); the cross-world switching rate \(\tau_J\) requires coupling assumptions or sharp bounds. The task is to distinguish corrections of mistaken beliefs from changes in preferences, social ties, or offer probabilities. A complementary structural target is the loss in expected job value attributable to inaccurate \(b_{ij}\), relative to truthful information with the same feasible offers. Random assignment identifies information effects only with consistency, appropriate interference restrictions, and measured attrition. At scale, employer responses and vacancy competition must be included. Determine when the intervention improves durable earnings and worker value, and characterize which belief dimensions and worker groups account for heterogeneity.

\textbf{Definitions and admissible scope.} Beliefs \(b_{ij}\) are probability measures over positive wages and an amenity space, and \(u_i(w,a)\) is a specified utility function. Fix finite job offers \(\mathcal O_i\), an outside option, switching costs \(k_{ij}\), and a tie-breaking rule. Under fixed offers, define \(j_i(b)\) as the maximizer of expected utility net of switching cost and define belief regret as the truthful expected utility net of switching cost at \(j_i(b^{\rm true})\) minus the corresponding net utility at \(j_i(b)\). Earnings \(Y_{iT}(z)\) include zeros during nonemployment and have finite means. For a finite employer set with marginal probabilities \(p_j^z\), randomization without interference identifies those marginals but not the joint law \((J_i(0),J_i(1))\); the sharp switching bounds are obtained by minimizing and maximizing \(\sum_{j\ne k}\gamma_{jk}\) over nonnegative couplings \(\gamma\) with row and column sums \(p^0,p^1\). Use dates \(t=0,\ldots,T\), fixed \(\delta\in(0,1]\), and \(Y_{iT}(z)=\sum_{t=0}^T\delta^t y_{it}(z)\), where \(y_{it}\) is real earnings. In the fixed-offer exercise, \(v_{ij}=\int u_i(w,a)\,db_{ij}^{\rm true}(w,a)\) and \(j_i(b)=\arg\max_{j\in\mathcal O_i\cup\{0\}}[\int u_i\,db_{ij}-k_{ij}]\) with the declared tie rule; switching cost at the outside option is specified as well. All expectations are taken over the fixed baseline worker measure and intervention uncertainty. For the finite-employer switching bound, include unemployment as employer code 0. Changes in offers or preferences are distinct extensions of this fixed-offer regret target, rather than automatically classified as belief corrections. Explicitly, let \(q_{ij}=v_{ij}-k_{ij}\). The fixed-offer belief regret is \(q_{i,j_i(b^{\rm true})}-q_{i,j_i(b)}\ge0\), with the same outside option and switching costs in both choices. A difference in gross values \(v\) would measure a different object and need not be nonnegative. Require every feasible offer and the outside option to have finite true and believed expected utility and finite switching cost; finite earnings alone does not guarantee that the utility-regret expression is defined.
\subsection*{Variants and assumption changes}
\begin{enumerate}
\item Nonpay uncertainty version (source-motivated): replace wage messages by randomized, validated messages about a particular amenity. Measure posterior beliefs, offers, switching, and net job value; do not interpret ties or preference changes as factual mistakes.
\item Saturation version (editorial): assign markets coverage \(q\in\{0,q_1,1\}\) and define \(Y_{iT}(z_i,q)\). Estimate direct and spillover effects separately, allowing vacancies and employer wages to change with coverage.
\end{enumerate}
\subsection*{References and source audit}
\textbf{Support and access.} Full manuscript checked. Nonpay uncertainty and social ties are stated directions; the source finds preferences important and does not establish pay misinformation as the dominant explanation. \textbf{Locator:} December 2025 manuscript, p. 41, final paragraph.
\begin{enumerate}\raggedright
\item\hypertarget{definitions:source:30007634:1}{} Sydnee Caldwell; Ingrid Haegele; Jörg Heining. 2025. \emph{Why Workers Stay: Pay, Beliefs, and Attachment}. December 2025 version, conclusion, printed p. 41, final paragraph.
\url{https://sydneec.github.io/Website/CHH_Monopsony.pdf}.
DOI: \url{https://doi.org/10.3386/w33445}.
\end{enumerate}

\subsection*{Integrated refinement: ECON-095}
{\sffamily\small\color{muted}Job-search assistance: displacement versus net job creation\par}
This refinement adopts the additional source conventions
(page~\pageref{audited:conventions}), subject to its local restrictions.
\textbf{Net job creation versus displacement of untreated competitors.}
For a fixed eligible labor market, define a full assignment law $\pi_s(z)$
at coverage $s$. Employment $E_i(z)$ depends on competitors' assignments and
endogenous vacancies. Identify
\[
\Delta J(s)=\mathbb E_{z\sim\pi_s}\!\left[\sum_iE_i(z)\right]
              -\mathbb E\left[\sum_iE_i(0)\right].
\]
Compare this net effect with gains to treated workers at the same horizon and
population. Saturation or market-level designs must identify losses to
competitors and vacancy creation; individual randomization alone need not.
Distinguish employer information, skills, wage responses and expansion of
positions. Include nonemployment, missing workers and linked labor markets in
the exposure and follow-up convention.
\subsection*{Supplied source audit for ECON-095}
Local [S1], [S2], etc. in this audit and the references immediately below
belong to ECON-095, independently of the original entry's reference numbering.
[S1] and [S2] explicitly support untreated displacement versus net job creation. The assignment-law formulation is editorial and makes the coverage experiment complete.
\subsection*{References for integrated refinement ECON-095}
{\footnotesize\raggedright
\par\noindent\textbf{[S1]} Stefano Caria, Kate Orkin, Alison Andrew, Robert Garlick, Rachel Heath, and Niharika Singh (2024). \emph{Barriers to Search and Hiring in Urban Labour Markets}. VoxDevLit 10(1), February 2024.\par \textit{Verified locator / source role:} Primary publisher page checked for review identity, authors, version and background. The specific published agenda is corroborated in [S2]; the exact formal target is editorial.\par \url{https://voxdev.org/voxdevlit/barriers-search-and-hiring-urban-labour-markets}\par\smallskip
\par\noindent\textbf{[S2]} Oliver Hanney / VoxDev (living compilation). \emph{Evidence gaps: Key unanswered questions in development economics}. Accessed 8 October 2026. The page currently covers 20 topics; the ``16-topics URL is a historical slug.\par \textit{Verified locator / source role:} Topic heading: Barriers to Search and Hiring in Urban Labour Markets. Supports the broad research agenda; this is not a verbatim quotation or an attribution of the displayed mathematical target.\par \url{https://voxdev.org/topic/evidence-gaps-key-unanswered-questions-across-16-topics-development-economics}\par\smallskip
}

\subsection*{Integrated refinement: ECON-096}
{\sffamily\small\color{muted}Firm productivity gains from improved worker matching\par}
This refinement adopts the additional source conventions
(page~\pageref{audited:conventions}), subject to its local restrictions.
\textbf{Firm-side productivity and social match surplus.}
Let $\mu(z)$ match workers to positions, allowing firms multiple positions and
unfilled vacancies. With a common-currency discounted gross output value less
training and real opportunity costs $v(q,d,k)$, identify
\[
\Delta V=\mathbb E\left[\sum_{(i,f)\in\mu(1)}v(q_{if},d_{if},k_{if})
       -\sum_{(i,f)\in\mu(0)}v(q_{if},d_{if},k_{if})\right]-C(z).
\]
Specify capacities, outside options, match duration, entry and intervention
dependence of output/costs. Include unmatched workers, rejected applicants and
disappearing vacancies; completed matches alone select the comparison.
Separate pure reassignment from vacancy or production expansion. Wages are
worker--firm transfers rather than an additional social resource cost unless
they explicitly proxy labor opportunity cost. Firm profit, earnings and social
match surplus are distinct targets.
\subsection*{Supplied source audit for ECON-096}
Local [S1], [S2], etc. in this audit and the references immediately below
belong to ECON-096, independently of the original entry's reference numbering.
[S1] and [S2] explicitly support firm-side productivity gains from matching interventions. The aggregate match-value function is editorial and requires the stated opportunity-cost and selection rules.
\subsection*{References for integrated refinement ECON-096}
{\footnotesize\raggedright
\par\noindent\textbf{[S1]} Stefano Caria, Kate Orkin, Alison Andrew, Robert Garlick, Rachel Heath, and Niharika Singh (2024). \emph{Barriers to Search and Hiring in Urban Labour Markets}. VoxDevLit 10(1), February 2024.\par \textit{Verified locator / source role:} Primary publisher page checked for review identity, authors, version and background. The specific published agenda is corroborated in [S2]; the exact formal target is editorial.\par \url{https://voxdev.org/voxdevlit/barriers-search-and-hiring-urban-labour-markets}\par\smallskip
\par\noindent\textbf{[S2]} Oliver Hanney / VoxDev (living compilation). \emph{Evidence gaps: Key unanswered questions in development economics}. Accessed 8 October 2026. The page currently covers 20 topics; the ``16-topics URL is a historical slug.\par \textit{Verified locator / source role:} Topic heading: Barriers to Search and Hiring in Urban Labour Markets. Supports the broad research agenda; this is not a verbatim quotation or an attribution of the displayed mathematical target.\par \url{https://voxdev.org/topic/evidence-gaps-key-unanswered-questions-across-16-topics-development-economics}\par\smallskip
}

\subsection*{Common conventions: Source mathematical conventions}

These conventions apply to each entry unless explicitly replaced.

\textbf{Models and identification.} A primitive model \(m\) specifies all probability laws, feasible choices, information, equilibrium and measurement rules used by its target. The admissible class \(\mathcal M\) is fixed independently of outcomes. If \(P_O(m)\) is its observable probability law and \(\tau(m)\) the target, its identified set at observable law \(P\) is
\[
 \mathcal I_\tau(P)=\{\tau(m):m\in\mathcal M,\ P_O(m)=P\}.
\]
Point identification means that this set is a singleton. An empty set signals incompatibility of data and assumptions. Coordinate bounds need not describe the joint set. Bounds are sharp only if every retained target is attainable or arbitrarily approachable by compatible models. Finite bounds require explicit restrictions on unobserved quantities; unrestricted confounding, hidden wealth, extrapolation or arbitrary equilibrium selection can yield uninformative sets.

\textbf{Causal interventions.} A potential outcome \(Y_i(p)\) is the outcome under a complete policy regime \(p\), including timing, financing, information and market spillovers. The target population and units are fixed before treatment. With interference, use the complete assignment vector or a justified exposure mapping; individual no-interference notation alone cannot identify an economy-wide intervention. A difference of mean potential outcomes does not require the cross-world joint distribution. Conditioning on post-treatment migration, survival, enrollment or employment changes the estimand and needs separate justification.

\textbf{Statistical statements.} An estimator, confidence set, forecast or test specifies a sampling law, model class and loss/coverage criterion. Uniform \((1-\alpha)\)-coverage means \(\inf_{m\in\mathcal M}\Pr_m\{\tau(m)\in C_n\}\geq1-\alpha\), or an explicitly asymptotic version. The whole parameter space trivially has coverage, so informativeness requires a stated width, risk or power objective. A forecasting improvement is relative to a named benchmark, information set and evaluation population.

\textbf{Equilibrium and optimization.} An equilibrium comprises feasible plans, optimality, beliefs where needed, budget constraints and all market/resource clearing. The primitives or a stated benchmark define each correspondence \(\mathcal E(p;m)\). Nonemptiness is assumed or proved, never inferred just from compact policy sets. A selected-equilibrium target uses a declared selection rule; a robust target quantifies over all permitted equilibria. Use suprema or infima unless attainment is established. An \(\varepsilon\)-optimal policy is feasible and within \(\varepsilon\) of the finite optimum value. Report an empty feasible set separately.

\textbf{Welfare and accounting.} Normative utility, interpersonal normalization, social weights, numeraire, discounting and the use of public receipts are primitives. Behavioral utility need not coincide with the welfare criterion. Household welfare includes ownership income and rents once; adding profits again to compensating variation generally double-counts them. A monetary equivalent is defined by an explicit transfer and price convention, with monotonicity and range conditions. Average effects, mechanism contributions, transfers and welfare losses are different targets. Infinite sums require convergence and dynamic feasible plans require terminal or no-Ponzi conditions.

\textbf{Source versions.} A working-paper date, revision date and journal publication year are distinguished. A DOI for a journal version is not automatically the DOI of an earlier manuscript. Locators refer to the stated version and printed pages where specified. A metadata-only check does not verify a claim about a full-text conclusion. Every entry's source subsection narrows the provenance claim accordingly.

\subsection*{Common conventions: Additional-source status and mathematical conventions}

\label{audited:conventions}

Of the 100 supplied ECON entries, 65 distinct problems appear here; 32 close
matches refine earlier entries, and three duplicate questions map to existing
entries. Appendix~\ref{app:audited-crosswalk} lists every destination.
The attachment's P/R/S/U distinctions and source audits are retained. Its
precise mathematical formalizations are editorial unless the individual audit
says otherwise. Candidate subsidy targets ECON-004--006 retain their U flags;
the resolved approval-core question ECON-007 updates OP-0202 above.
The following supplied conventions govern both this part and the attributed
ECON refinements elsewhere in the catalogue, subject to local restrictions.

\textbf{Parameterized questions.} A problem family lists inputs that must be fixed before an instance is evaluated: population, horizon, policies, information, data law, model class and welfare weights. These are required choices, not quantities the citations are assumed to specify. Undefined applications should not be treated as identified estimands.

\textbf{Indices and integrability.} Sets of agents, firms and regions are finite unless a population measure is explicitly used. \(i,j\) index units, \(r\) regions, \(t\) dates and \(h\) horizons. \(\mathbb E\) and \(\Pr\) refer to the declared population and law; \(\mathbf1\{A\}\) is an indicator. Every displayed expectation and discounted sum needs existence in the stated domain. Infinite horizons use a discount in \((0,1)\) and a convergence condition; finite horizons may use discount one.

\textbf{Optimization and existence.} A displayed \(\max\), \(\min\), or \(\arg\max\) is conditional on attainment. For finite feasible menus this is immediate; general spaces need continuity and compactness/coercivity or another existence argument. Without attainment, the target is the supremum/infimum and arbitrarily near-optimal feasible choices. \(\inf\varnothing=+\infty\). Empty feasibility and an unidentified target are different issues.

\textbf{Potential outcomes.} \(Y_i(z)\) is the outcome under a specified intervention, not the observed conditional mean among people choosing \(z\). In binary comparisons \(\Delta Y_i=Y_i(1)-Y_i(0)\). Specify assignment, treatment versions, compliance, information and observation horizon. Interference is allowed under a full policy vector; compressed exposure notation needs a justified mapping. No interference, consistency and overlap are substantive assumptions, not automatic consequences of notation.

\textbf{Populations and observation.} Fix baseline eligibility and population weights. Conditioning on post-policy employment, survival, relocation or program uptake can change the population being compared. Death is not ordinary loss to follow-up; outcomes truncated by death need a composite convention, joint survival target, or explicitly defined survivor stratum. Fertility-changing interventions need a population functional rather than an unexplained fixed descendant cohort. Measurement and missingness models must be supplied.

\textbf{Identification.} A structural model \(m\in\mathcal M\) implies observable law \(P_m\) and target \(\theta(m)\). Identification means
\[
 P_{m_1}=P_{m_2}\ \Longrightarrow\ \theta(m_1)=\theta(m_2)
 \quad\text{for all }m_1,m_2\in\mathcal M .
\]
Without identification, the target is
\[
 \Theta(P;\mathcal M)=\{\theta(m):m\in\mathcal M,\ P_m=P\}.
\]
Writing \(\theta(P)\) before establishing identification can hide incompatible latent models. Confidence sets additionally need a sampling law, confidence level, asymptotic sequence and uniformity class. Point identification, coverage and informative precision are separate properties.

\textbf{Mediation.} Total effects of randomized assignment are distinct from cross-world natural indirect effects. Belief/effort-channel effects require a meaningful mediator intervention and additional measurement, exclusion or mediation assumptions. Randomization of the initial policy alone does not supply those assumptions. Report bounds or interventional alternatives when the stronger target is unsupported.

\textbf{Equilibrium.} A counterfactual \(W(z)\), \(p(z)\), or \(\mu_t^z\) requires individual optimization, feasibility, government/resource budgets, market clearing and state-transition laws. State equilibrium existence and selection, or report ranges over compatible equilibria. Initial-state integration includes subsequent endogenous transitions; current-state integration must use the policy-dependent distribution. Reduced-form invariance after policy is not assumed.

\textbf{Welfare accounts.} Scores, utility, health, dollars and ecological quantities require explicit conversion before addition. Fees, taxes, subsidies, wages and extortion payments are transfers between parties; distributional weights can make them welfare relevant, but their face value is not automatically a resource loss. State ownership, geographic boundaries, foreign-income weights and fiscal finance. Do not add profits to household welfare already including dividends, or count land capitalization and the underlying benefit twice. A benefit-cost expression must distinguish gross benefits from net welfare and subtract every real cost exactly once.

\textbf{Derivatives and risk.} Log derivatives require positive arguments; local responses require admissible perturbations and specified equilibrium branches. Differentiation under expectations needs regularity/domination, and boundaries may allow only one-sided derivatives. Risk uses a specified law; ambiguity uses a declared nonempty law set on a common history space. Adaptive policies must be nonanticipative. A robust criterion is an editorial choice unless attributed explicitly.

\textbf{Scope overlaps.} Allocation, collective choice, contracts, household bargaining and coordinated policy overlap economics with optimization and game theory. They are retained because they appeared in the original catalogue; no claim of a strictly game-theory-free selection is made.
% END ECONOMICS STATEMENT
\end{document}
